\documentclass[10pt,a4paper]{article}
\usepackage{amsmath,amssymb,amsthm,graphicx}
\usepackage{xfrac}
\newenvironment{enumeratei}
{\bgroup\def\theenumi{\roman{enumi}}\def\theenumii{\arabic{enumii}}\begin{enumerate}}
{\end{enumerate}\egroup}

\hyphenation{addi-tion-al belong belongs denote exchan-ge exchan-ged exchan-ges}

\newcommand{\psfrac}[2]{\sfrac{(#1)}{#2}}

\newcommand\mto{\mathchoice{\longmapsto}{\mapsto}{\mapsto}{\mapsto}}
\let\oldbigsqcup\bigsqcup
\renewcommand{\bigsqcup}{\mathop{\textstyle\oldbigsqcup}\displaylimits}

\let\epsilon\varepsilon

\DeclareMathOperator{\Ess}{Ess}
\DeclareMathOperator{\Leb}{Leb}

\usepackage{tikz}
\usetikzlibrary{plotmarks}

\theoremstyle{plain}
\newtheorem{theorem}{Theorem}[section]
\newtheorem{lemma}[theorem]{Lemma}
\newtheorem{proposition}[theorem]{Proposition}
\newtheorem{corollary}[theorem]{Corollary}
\newtheorem{claim}{Claim}

\theoremstyle{definition}
\newtheorem{remark}[theorem]{Remark}
\newtheorem{example}[theorem]{Example}

\newcommand{\A}{\mathcal{A}}
\newcommand{\B}{\mathcal{B}}
\newcommand{\I}{\mathcal{I}}

\newcommand{\Z}{\mathbb{Z}}
\newcommand{\N}{\mathbb{N}}
\newcommand{\R}{\mathbb{R}}
\title{Iterates of the endpoints of the exchanged intervals produce symmetric intervals}
\author{Przemyslaw Berk, Frank Trujillo}
\date{doi:10.5802/jep.302; CC BY4.0}
\begin{document}
\maketitle
\noindent\textit{Editorial note.} The counted statement is the author's labelled lemma environment at source0.tex lines 693-704; its proof environment opens at line 706, the line after the blank separator, so the pairing is by adjacency and the retained proof range 706-724 is the whole environment.\par\medskip
\noindent\textit{Editorial note.} Three numbered identities are retained as context because the counted proof combines them: eq: standardconjugacy (source lines 463-465), eq:generalized\_conjugacy (source lines 469-472) and eq:sym\_identity\_endpoints (source lines 475-478).\par\medskip
\noindent\textit{Editorial note.} This source is new to the material-authoring work and its preamble is built from the author's own macro and theorem block (source lines 8-45) rather than from a prepared adapter.\par\medskip
\noindent\textit{Editorial note.} This source has no label-number dump and the retained counter values are left empty, so the printed numbering of the card is the wrapper's default; the identity of the counted item is its author label.\par\medskip
\noindent\textit{Editorial scope.} One original result of the pinned publisher article, counted once: iterates of the endpoints of the exchanged intervals span symmetric subintervals (author label lem: sym\_endpoints).\par\medskip
\noindent\textit{Allowed standard prerequisites (editorial).}\begin{itemize}
\item Connection time of an endpoint. M(alpha) is the first time the orbit of the left endpoint of the exchanged interval I\_alpha meets the set of discontinuities.. The hypothesis of the statement bounds the exponent m by it.
\item Conjugacy identities. The reflection conjugates the transformation to its inverse and sends the endpoint of an exchanged interval to the endpoint of the next exchanged interval.. They are the two identities combined in the computation.
\item Symmetric interval. A subinterval is symmetric when the reflection of the whole interval maps it to itself by exchanging its endpoints.. It is the property concluded in the second half of the statement.
\end{itemize}\medskip
\begin{equation}\label{eq: standardconjugacy}
\mathcal{I}_I\circ T(x)=T^{-1}\circ\mathcal{I}_I(x) \quad \text{if } x \neq \partial I_{\alpha} \text{ for any } \alpha \in \A.
\end{equation}

\begin{multline}
\label{eq:generalized_conjugacy}
\mathcal{I}_I\circ T^n(x)=T^{-n}\circ\mathcal{I}_I(x) \quad \text{if } x \neq T^{-i}(\partial I_\alpha) \text{ for any } \alpha \in \A \text{ and}\\ \text{any $i$ such that } 0 \leq i < n.
\end{multline}

\begin{equation}
\label{eq:sym_identity_endpoints}
\I_I\circ T(\partial I_{\alpha})=\partial I_{\hat\alpha},\ \text{where}\ \pi_0(\hat\alpha)=\pi_0(\alpha)+1,
\end{equation}

\begin{lemma}\label{lem: sym_endpoints}
{Let $T : I \to I$ be a symmetric IET.}
Let $\alpha\in\mathcal A\setminus\{\pi_0^{-1}(1)\}$ and $m< M(\alpha)$. Then
\begin{equation}\label{eq: connections}
\I_I\circ T(T^{-m}(\partial I_{\alpha}))={\I_I\circ (T^{-m+1}(\partial I_{\alpha}))}=T^m(\partial I_{\hat\alpha}),
\end{equation}
where $\pi_0(\hat\alpha)=\pi_0(\alpha)-1$.

In particular, the left-closed right-open interval with endpoints $T^{-m}(\partial I_{\alpha})$ and $T^m(\partial I_{\hat\alpha})$ is {$\beta$-symmetric for some $\beta\in\A$, while the left-closed right-open interval with endpoints $T^{-m+1}(\partial I_{\alpha})$ and $T^m(\partial I_{\hat\alpha})$ is $\sfrac{1}{2}$-symmetric}.

Moreover, $M(\alpha)= N(\hat\alpha)$.
\end{lemma}

\begin{proof}
First, notice that \eqref{eq: connections} is equivalent to \eqref{eq:sym_identity_endpoints} if $m = 1$. This implies that $\I_I(\partial I_\alpha) = T(\partial I_{\hat\alpha})$.
In the following, we~assume $m > 1$. If $m < M(\alpha)$, by~\eqref{eq:generalized_conjugacy},
\[
\I\circ T(T^{-m}(\partial I_{\alpha}))= \I \circ T^{-m + 1}(\partial I_\alpha) = T^{m - 1} \circ \I (\partial I_\alpha) = T^m(\partial I_{\hat \alpha}).
\]
This proves \eqref{eq: connections}, which easily implies that the symmetry properties of the intervals in the statement.

Assume now that $m=M(\alpha)$ and $T^{-m}(\partial I_{\alpha})=\partial I_{\beta}$ for some $\beta\in \A$. Since $M(\alpha)$ is the minimal number that makes a connection, by~\eqref{eq: connections} we have
\[
\I\circ T(T(\partial I_{\beta}))=\I\circ T(T^{-{m+1}}\partial I_{\alpha})=T^{m-1}(\partial I_{\hat\alpha}).
\]
Moreover, noticing that $T(\partial I_\beta) \notin \{\partial I_\alpha\}_{\alpha \in \A}$ since $m > 1$, by~\eqref{eq: standardconjugacy}
\[
\I\circ T(T(\partial I_{\beta}))=T^{-1}\circ\I\circ T(\partial I_{\beta})
=T^{-1}(\partial I_{\hat\beta}),
\]
where $\pi_0(\hat\beta)=\pi_0(\beta)+1$. Thus we get $T^{m}(\partial I_{\hat\alpha})=\partial I_{\hat\beta}$. This proves $N(\hat\alpha)\le M(\alpha)$. The opposite inequality is proved analogously, {by exchanging the roles of $T$ and $T^{-1}$.}
\end{proof}

\end{document}
